\section{A fronteira estocástica: o que foi calculado e por que o
cálculo não serve como medida}
\label{ap:sfa}

% O prefixo segue \thesection em vez da letra fixa: assim a numeração dos
% flutuantes acompanha a posição do apêndice se a ordem mudar, do mesmo
% modo que as referências cruzadas simbólicas já acompanham.
\setcounter{figure}{0}
\renewcommand{\thefigure}{\thesection.\arabic{figure}}
\setcounter{table}{0}
\renewcommand{\thetable}{\thesection.\arabic{table}}

Este apêndice existe por uma razão de honestidade de registro. A
Seção~\ref{sec:fase2} anuncia que a eficiência deste documento é medida
por um método só, a análise envoltória de dados, e que um segundo método
foi calculado e \textbf{deu errado}. Aqui está o quê, exatamente, deu
errado. Quem quiser apenas os resultados de eficiência não precisa desta
página: nenhum número dela entra em conclusão alguma do documento.

\subsection*{O que se tentou medir, e por qual caminho}

Existem duas famílias de método para medir eficiência hospitalar, e a
pesquisa tentou as duas. A que sustenta os resultados, a
\textbf{análise envoltória de dados}, não supõe fórmula nenhuma: olha o
que os hospitais de fato conseguiram fazer com o que tinham, traça a
fronteira de melhores práticas pelos melhores casos observados e mede a
distância de cada um até ela. A outra, a \textbf{fronteira estocástica},
faz o contrário: parte de uma equação escrita de antemão para
transformar recursos em atendimentos e, com ela, tenta separar o que é
ineficiência de gestão do que é apenas azar do ano.

O cálculo foi feito por inteiro, e não parou por falta de ajuste
aparente. A forma funcional adotada foi a \textbf{Translog}, uma forma
matemática mais flexível que a alternativa clássica, a Cobb-Douglas, por
não obrigar os insumos a se combinarem sempre na mesma proporção; ela
foi preferida com folga larga (razão de verossimilhança de \textbf{567,0}
com 6 graus de liberdade, $p$ da ordem de $10^{-119}$), e
\textbf{95,5\%} da variância composta foi atribuída à ineficiência. O
resultado, além disso, \textbf{apontou na mesma direção do DEA}: as OSS
na ponta de cima, o mesmo ordenamento geral entre as redes.

É por isso que o problema precisa ser dito com precisão. O método não
falhou em produzir números, nem produziu números que contrariassem o
resto do documento. Ele produziu números que \textbf{não medem o que o
nome diz que medem}, por três defeitos independentes.

\subsection*{Defeito 1: o modelo trava exatamente no grupo de interesse}

\textbf{Este é o defeito que inutiliza o método para esta pesquisa.} A
fronteira devolveu, para \emph{todas} as OSS, praticamente o mesmo
valor: uma eficiência técnica que fica entre \textbf{0,9993 e 0,9994} em
todo o painel, variando menos de dois centésimos de milésimo entre o
primeiro e o último hospital da categoria, ao longo de onze anos.

Isso não é um resultado dizendo que as OSS são parecidas entre si; é o
modelo travando. Um instrumento que dá a mesma nota a todos os hospitais
de um grupo não está medindo aquele grupo, e as OSS são justamente a
categoria que esta pesquisa existe para comparar.

A consequência aparece na segunda equação do modelo, a que explica
\emph{quanta} ineficiência cada hospital carrega: as OSS ficam tão
coladas à fronteira que o coeficiente delas ali \textbf{perde
identificação}, que é o análogo da separação perfeita já discutida na
Seção~\ref{sec:estrategia}. A assinatura do problema é a instabilidade
do ponto, não a sua imprecisão aparente: nas quatro amostras e nas três
sensibilidades o coeficiente oscila entre $-772$ e $-1.971$, com
erro-padrão entre 400 e 1.896. Um parâmetro identificado não varia por
fator de 2,6 entre subamostras da mesma base.\footnote{Numa das amostras
(sem o biênio pandêmico) o valor sai com $p$ de $8{,}3\times10^{-7}$.
Isso \emph{não} é precisão: sob separação a verossimilhança fica
praticamente plana numa região distante de zero, e o intervalo de Wald
calculado ali não tem interpretação. Um $p$ pequeno sobre um parâmetro
degenerado mede o formato da aproximação, não o dado.} O defeito
reaparece em todas as variantes testadas e não tem conserto dentro deste
método.

\subsection*{Defeito 2: a fronteira não se comporta como função de produção}

Uma função de produção precisa ser \textbf{monotônica}: acrescentar
insumo não pode reduzir o produto possível. Na fronteira estimada, o
insumo de UTI viola essa condição em \textbf{94,9\% das observações}
(97,8\% entre as que têm leito de UTI habilitado, tenham-no usado ou não
no ano; 96,9\% se o corte for ter registrado diária de UTI), contra
99,7\% de respeito para leitos e 99,6\% para valor. É herança do
logaritmo de um mais UTI, com 19\% de zeros, e da colinearidade alta com
leitos totais e valor.

\textbf{A violação não se concentra em categoria nenhuma}, o que impede
tratá-la como problema de uma linha só: vai de 93,6\% no Filantrópico a
\textbf{99,3\% na administração Direta}, que é justamente a referência de
todos os contrastes. Cinco vírgula sete pontos de amplitude sobre uma
base de 93,6\% é defeito geral da forma funcional. \textbf{A fronteira
inteira está desqualificada como medida}, não uma parte dela.

Isto tem uma consequência que costuma passar despercebida:
significância mede a precisão de uma estimativa, não a validade da
estrutura de onde ela sai. Coeficiente preciso extraído de fronteira
mal-comportada é preciso sobre um objeto que não é o que o nome diz.

\subsection*{Defeito 3: um erro de montagem do insumo, que ajuda a
entender os outros dois}

A lista de recursos entregue ao modelo continha o total de leitos do
hospital \emph{e}, separadamente, os leitos de UTI, que já estavam
contados dentro do total. É como perguntar quanto custa um carro e
depois, na mesma conta, quanto custa o motor dele, somando as duas
coisas. O modelo então precisa responder o que acontece quando se
acrescenta um leito de UTI \emph{sem} acrescentar um leito ao hospital,
que é uma pergunta sem sentido prático, e devolve uma resposta sem
sentido: mais UTI aparecendo como se reduzisse a capacidade.

Refeito o cálculo sem essa duplicação, o problema desaparece e a ordem
dos hospitais quase não muda, o que mostra que o defeito estava na
montagem e não na realidade hospitalar que se queria descrever. O
diagnóstico importa por si: sem ele, a violação de monotonicidade do
Defeito 2 seria lida como característica do parque hospitalar, e não
como artefato da lista de insumos.

\subsection*{A tabela que se preservou, e o estatuto dela}

A equação de ineficiência é reproduzida abaixo porque omiti-la seria
pior: ela contém o único resultado desta fase que o documento chegou a
não mostrar, e resultado que desaparece depois de encontrado é omissão.
Ela vem, porém, com a advertência dos parágrafos anteriores colada a
cada uma das suas linhas.

\begin{table}[H]
\centering
\caption{Equação de ineficiência da fronteira estocástica (Battese e
Coelli 1995). Referência: administração Direta, porte médio.
Coeficiente negativo indica \emph{menos} ineficiência.
\textbf{Tabela tarjada: a monotonicidade da fronteira que produziu estes
coeficientes está violada em 94,9\% das observações, de 93,6\% no
Filantrópico a 99,3\% na própria Direta. Nenhum valor desta tabela é
interpretável como magnitude nem entra em conclusão deste documento; ela
é diagnóstico da inadequação da forma Translog.}}
\label{tab:equacaoinef}
\small
\begin{tabular}{lrrr}
\toprule
Termo & Coeficiente & Erro-padrão & $p$ \\
\midrule
Público Municipal        & $-5{,}454$ & 1,020 & $9{,}0\times10^{-8}$ \\
Filantrópico             & $-0{,}244$ & 0,163 & 0,134 \\
Privado                  & $+2{,}474$ & 0,396 & $4{,}2\times10^{-10}$ \\
Complexidade (padronizada) & $-2{,}159$ & 0,325 & $3{,}2\times10^{-11}$ \\
Grande Porte             & $+2{,}470$ & 0,333 & $1{,}3\times10^{-13}$ \\
Especial                 & $+7{,}401$ & 0,917 & $7{,}0\times10^{-16}$ \\
Tendência linear         & $-0{,}010$ & 0,023 & 0,668 \\
\bottomrule
\end{tabular}
\end{table}

\textbf{Lida como diagnóstico e não como medida}, a linha que mais
importa é a do Público Municipal ($-5{,}454$, $p$ de
$9{,}0\times10^{-8}$): precisa, e apontando o Público Municipal como
menos ineficiente que a administração Direta. \textbf{O DEA, que é o
instrumento adequado, não confirma essa ordem, e tampouco a refuta.} Nas
medianas corrigidas de viés da Seção~\ref{sec:fase2} a Direta (0,634)
fica à frente do Município (0,632); nas médias o Município fica à frente
(0,644 contra 0,620). Ano a ano, no escore corrigido de viés, o Município
supera a Direta em seis dos onze anos, e os seis são de 2020 a 2025. No
estrato de alta complexidade da Seção~\ref{sec:segmentacao}, que é onde a
comparação entre redes é viável, o Município fica à frente; no estrato
intermediário, atrás. \textbf{A ordem entre estas duas redes depende de
qual escore e de qual estatística se escolhe, e por isso o documento não
a afirma em direção nenhuma.} O que há é um sinal recorrente e sem
corroboração por instrumento adequado, o que é diferente tanto de achado
quanto de nada.

As demais linhas têm o mesmo estatuto. O Filantrópico não se distingue da
Direta ($p$ de 0,134). O porte carrega o maior coeficiente da tabela
(Especial, $+7{,}401$, quase três vezes o do Grande Porte) enquanto a
complexidade anda no sentido oposto ($-2{,}159$), e não há tendência
temporal ($p$ de 0,668). Lidos como \emph{ordens de grandeza relativas
dentro de uma fronteira que se sabe mal especificada}, esses sinais
sugerem que a estrutura física do hospital move a ineficiência medida
mais do que a rede que o administra. Lidos como magnitudes, não dizem
nada.

\subsection*{Concordar com o DEA não conserta nada}

\comoler{quanto mais a nuvem sobe da esquerda para a direita, mais os
dois procedimentos se parecem ao ordenar os hospitais. A faixa colada no
topo é a degeneração das OSS do Defeito 1: o modelo paramétrico as
coloca todas na fronteira. Mesmo descontando essa faixa, a semelhança
permanece, e mesmo assim ela não é confirmação: o eixo vertical vem de
uma fronteira cuja monotonicidade está violada em 94,9\% das
observações.}

\begin{figure}[H]
\centering
\includegraphics[width=.58\textwidth]{figS02_dispersao_sfa_dea.png}
\caption{Semelhança entre os dois procedimentos de eficiência: cada ponto
é um hospital num ano, posicionado pelo escore DEA (horizontal) e pela
eficiência da fronteira estocástica (vertical). O eixo vertical é
diagnóstico, não medida: nenhum valor dele é citável.}
\label{fig:convergencia}
\end{figure}

\comofeito{gráfico de dispersão. Cada ponto é um hospital em um ano,
posicionado pelo \textbf{escore de eficiência DEA} no eixo horizontal e
pela \textbf{eficiência técnica da fronteira estocástica} no vertical
(ambos entre 0 e 1). Quanto mais a nuvem sobe da esquerda para a
direita, mais os dois procedimentos, de naturezas opostas, concordam
sobre quem é eficiente.}

A correlação de ordem entre os dois conjuntos de escores é de
\textbf{0,62} (0,63 excluindo as OSS, de modo que a semelhança não
depende do grupo problemático). Ela não é corroboração de um método pelo
outro: os dois consomem exatamente a mesma informação, e um deles não
descreve uma função de produção. A semelhança tampouco cresce de forma
gradual, como cresceria se refletisse aprendizado do modelo: fica
achatada entre 0,50 e 0,55 de 2015 a 2019, salta para a faixa de 0,66 a
0,74 já em 2020 e se mantém nesse patamar até 2025, um degrau que
coincide com o início da pandemia.

\subsection*{A regra, enunciada uma vez para todo o documento}

\textbf{Nenhum coeficiente e nenhum nível da fronteira estocástica é
citável, em nenhuma categoria de gestão, em nenhuma seção.} O que dela se
aproveita é um achado de método, e não de gestão: a violação de
monotonicidade em 94,9\% das observações diz que a forma Translog
\textbf{não descreve bem o modo como estes hospitais transformam insumo
em produto}. \textbf{Todas as medidas de eficiência deste documento vêm
da análise envoltória de dados}, e é só sobre ela que as conclusões de
eficiência se apoiam.
